% hw3-predicate-logic.tex

\newpage
\section{作业 3：谓词逻辑}

\begin{problem}[一阶谓词逻辑: 形式化描述与推理]
  请使用一阶谓词逻辑公式描述以下两个定义,
  并从逻辑推理的角度说明这两种定义之间是否有强弱之分
  （可以使用``重言式''进行公式转换）。
  \noindent {(要求: 需给出关键的推理步骤或理由)}

  A function $f$ from $\mathbb{R}$ to $\mathbb{R}$ is called
  \begin{enumerate}[(1)]
    \setlength{\itemsep}{10pt}
    \item \blue{\it pointwise continuous} (连续的) if
      for every $x \in \mathbb{R}$
      and every real number $\epsilon > 0$,
      there exists real $\delta > 0$ such that
      for every $y \in \mathbb{R}$ with $|x - y| < \delta$,
      we have that $|f(x) -  f(y)|< \epsilon$.
    \item \blue{\it uniformly continuous} (一致连续的) if
      for every real number $\epsilon > 0$,
      there exists real $\delta > 0$ such that
      for every $x, y \in \mathbb{R}$ with $|x - y| < \delta$,
      we have that $|f(x) -  f(y)|< \epsilon$.
  \end{enumerate}
\end{problem}

\begin{solution}
  取论域为 $\mathbb{R}$。
	形式化表示如下:
  \setcounter{equation}{0}
  \begin{enumerate}[(1)]
    \item
      \begin{align}
        \forall x \in \mathbb{R}.\; \forall \epsilon \in \mathbb{R}^{+}.\;
          \exists \delta \in \mathbb{R}^{+}.\;
          \forall y \in \mathbb{R}.\; |x - y| < \delta \to |f(x) - f(y)| < \epsilon.
        \label{eq:cont}
      \end{align}
    \item
      \begin{align}
        \forall \epsilon \in \mathbb{R}^{+}.\; \exists \delta \in \mathbb{R}^{+}.\;
          \forall x \in \mathbb{R}.\; \forall y \in \mathbb{R}.\; |x - y| < \delta \to |f(x) - f(y)| < \epsilon.
        \label{eq:uniform-cont}
      \end{align}
  \end{enumerate}

  下面从逻辑的角度讨论 (\ref{eq:cont}) 与 (\ref{eq:uniform-cont}) 的强弱关系。

  首先, 根据重言式
  \[
    \forall x.\; \forall y.\; \alpha \leftrightarrow \forall y.\; \forall x.\; \alpha,
  \]
  (\ref{eq:cont}) 等价于
  \begin{align*}
    \forall \epsilon \in \mathbb{R}^{+}.\; \forall x \in \mathbb{R}.\;
      \exists \delta \in \mathbb{R}^{+}.\;
      \forall y \in \mathbb{R}.\; |x - y| < \delta \to |f(x) - f(y)| < \epsilon.
  \end{align*}

  其次, 根据重言式
  \[
    \exists x.\; \forall y.\; \alpha \to \forall y.\; \exists x.\; \alpha,
  \]
  可知 (\ref{eq:uniform-cont}) \red{不弱于} (\ref{eq:cont})。

  最后, 由于
  \[
    \forall y.\; \exists x.\; \alpha \nvdash \exists x.\; \forall y.\; \alpha,
  \]
  可知 (\ref{eq:uniform-cont}) (严格)强于 (\ref{eq:cont})。

  作为例子, $f(x) = x^{2}$ 是连续函数, 但不是一致连续函数。
\end{solution}